% ROCF 电子测评系统 — 设计与开发文档
% tectonic rocf_doc.tex
\documentclass[12pt,a4paper]{ctexart}

% === 页面布局 ===
\usepackage[top=2.5cm, bottom=2.5cm, left=2.8cm, right=2.5cm]{geometry}

% === 字体 ===
\setCJKmainfont{PingFang SC}
\setmonofont{Menlo}[Scale=0.9]

% === 常用宏包 ===
\usepackage{graphicx}
\usepackage{booktabs}
\usepackage{hyperref}
\usepackage{xcolor}
\usepackage{fancyhdr}
\usepackage{enumitem}
\usepackage{listings}
\usepackage{titlesec}
\usepackage{longtable}
\usepackage{array}
\usepackage{tabularx}
\setlength{\headheight}{14.5pt}

% === 超链接 ===
\hypersetup{
  colorlinks=true,
  linkcolor=blue!60!black,
  urlcolor=cyan!60!black,
  citecolor=green!50!black,
  linktoc=none
}

% === 代码块 ===
\lstset{
  basicstyle=\ttfamily\footnotesize,
  backgroundcolor=\color{gray!8},
  frame=single,
  framesep=6pt,
  rulecolor=\color{gray!40},
  breaklines=true,
  breakatwhitespace=false,
  language=Python,
  keywordstyle=\color{blue!60!black},
  commentstyle=\color{gray!50},
  stringstyle=\color{orange!70!black},
  showstringspaces=false,
  numbers=left,
  numberstyle=\tiny\color{gray},
  numbersep=8pt,
  tabsize=4,
  xleftmargin=16pt,
  framexleftmargin=16pt
}

% === 章节标题 ===
\titleformat{\section}{\Large\bfseries}{\thesection}{1em}{}
\titleformat{\subsection}{\large\bfseries}{\thesubsection}{1em}{}
\titleformat{\subsubsection}{\normalsize\bfseries}{\thesubsubsection}{1em}{}

% === 页眉页脚 ===
\pagestyle{fancy}
\fancyhf{}
\fancyhead[L]{ROCF 电子测评系统 — 设计与开发文档}
\fancyhead[R]{v0.1}
\fancyfoot[C]{\thepage}
\renewcommand{\headrulewidth}{0.4pt}

\begin{document}

% ============================================================================
% 标题页
% ============================================================================
\begin{titlepage}
\centering
\vspace*{6cm}

{\Huge\bfseries ROCF 电子测评系统\\[6pt] 设计与开发文档\par}

\vspace{2cm}

{\LARGE 版本：v0.1\par}

\vspace{1cm}

{\Large 日期：2026-07-27\par}

\vspace{1cm}

{\normalsize 广西幼儿师范高等专科学校艺术设计专业团队\\顽主（天津）教育科技有限责任公司\\联合开发\par}

\vfill

{\footnotesize 广西·南宁\par}
\end{titlepage}

% ============================================================================
% 目录
% ============================================================================
\tableofcontents
\newpage

% ============================================================================
% 正文
% ============================================================================

\section{项目背景与文献基础}

\subsection{ROCF 测验概述}

Rey-Osterrieth 复杂图形测验（ROCF）由 André Rey 于 1941 年首创，1944 年由 Paul-Alexandre Osterrieth 标准化，是神经心理学领域使用最广泛的视空间能力与记忆评估工具之一。其核心范式为"临摹→即时回忆→延迟回忆"，能够在单次施测中同时评估视空间建构能力、执行功能（规划与组织）及情景记忆，已翻译并适用于全球 40 余个国家。

\subsection{关键文献发现（本次开发依据）}

\begin{longtable}{p{4cm} p{5cm} p{5cm}}
\toprule
\textbf{维度} & \textbf{核心发现} & \textbf{对系统设计的指导} \\
\midrule
评分标准 & Osterrieth 18 单元 0--2 分制是国际通用基准（0--36 分） & 采用 18 单元评分作为默认系统，预留 Meyers 扩展接口 \\
\midrule
认知机制 & 临摹依赖额-顶叶执行网络，回忆依赖海马-前额叶回路 & 分别记录临摹与回忆阶段的笔序、潜伏期、组织策略 \\
\midrule
组织策略 & 整体式策略者回忆显著优于零散式（Savage 等, 2000） & 捕获绘制顺序，自动判定组织策略类型 \\
\midrule
储存/提取分离 & 再认测验可分离储存缺陷 vs 提取缺陷（Meyers \& Meyers） & 纳入再认测验模块 \\
\midrule
数字化前沿 & CNN 自动评分 MAE=0.95--1.97，平板施测可捕获人类不可见的笔压/笔序特征 & 完整记录笔画数据，为未来 AI 评分预留数据接口 \\
\midrule
跨文化常模 & 教育解释回忆方差 21--41\%，中国人群回忆得分低于西方常模（Zhang 等, 2018） & 内置中文常模参考，支持教育水平校正 \\
\bottomrule
\caption{关键文献发现汇总}
\end{longtable}

\section{系统设计目标}

\subsection{核心目标}

\begin{enumerate}[leftmargin=*]
  \item \textbf{完整复现标准 ROCF 三阶段施测流程}：临摹 → 干扰任务 → 延迟回忆
  \item \textbf{精确数字化记录}：捕获每一笔的坐标、时间戳、压力（平板端）、顺序
  \item \textbf{多评分体系兼容}：默认 Osterrieth 18 单元，可扩展 BQSS / Meyers 等
  \item \textbf{跨平台部署}：macOS、Windows、Linux 原生 GUI
  \item \textbf{数据标准化输出}：JSON 结构化数据 + PNG 截图，便于研究级分析
\end{enumerate}

\subsection{非功能需求}

\begin{table}[h]
\centering
\begin{tabular}{ll}
\toprule
\textbf{指标} & \textbf{目标} \\
\midrule
启动时间     & $<$ 2 秒 \\
内存占用     & $<$ 150 MB \\
数据保存延迟 & $<$ 100 ms（每笔结束后异步写入） \\
支持操作系统 & macOS 12+, Windows 10+, Ubuntu 20.04+ \\
离线可用     & 是，无需网络连接 \\
\bottomrule
\end{tabular}
\caption{非功能需求指标}
\end{table}

\newpage
\section{系统架构}

\subsection{总体架构}

\begin{figure}[htbp]
\centering
\includegraphics[width=\textwidth]{../output/rocf_architecture.pdf}
\caption{ROCF 电子测评系统架构图}
\label{fig:architecture}
\end{figure}

\newpage
\subsection{技术栈}

\begin{table}[h]
\centering
\small
\begin{tabularx}{\textwidth}{l>{\raggedright\arraybackslash}p{4.2cm}>{\raggedright\arraybackslash}X}
\toprule
\textbf{层级} & \textbf{技术选型} & \textbf{理由} \\
\midrule
GUI 框架   & PySide6 (Qt 6)  & 原生 macOS 外观，跨平台，低资源占用 \\
图形渲染   & QPainter + QPainterPath & 硬件加速，支持抗锯齿与贝塞尔曲线 \\
数据存储   & JSON (本地文件)  & 人类可读，研究工具链兼容（Python/R/SPSS） \\
截图       & QWidget.grab() → PNG & 标准格式，可直接用于评分分析 \\
打包分发   & py2app (macOS) / PyInstaller (Win/Linux) & 独立 .app / .exe，无需安装 Python \\
字体       & PingFang SC / Microsoft YaHei / Noto Sans CJK & 跨平台中文字体自动适配 \\
\bottomrule
\end{tabularx}
\caption{技术栈选型}
\end{table}

\subsection{页面状态机}

\begin{figure}[htbp]
  \centering
  \includegraphics[width=0.90\textwidth]{../output/rocf_statemachine.pdf}
  \caption{ROCF 电子测评系统页面状态机}
  \label{fig:statemachine}
\end{figure}

\newpage
\subsection{数据模型}

\begin{lstlisting}[frame=single, basicstyle=\ttfamily\footnotesize, backgroundcolor=\color{gray!8}]
{
  "subject": {
    "id": "DEMO001",
    "age": "30",
    "gender": "男",
    "hand": "右利手"
  },
  "timestamp": "20260727_182900",
  "copy": {
    "phase": "copy",
    "strokes": [
      {
        "index": 0,
        "tool": "thin",
        "color": "#000000",
        "points": [
          {"x": 350, "y": 250, "t": 0.0},
          {"x": 355, "y": 251, "t": 0.016}
        ]
      }
    ],
    "total_time_ms": 45000,
    "stroke_count": 8
  },
  "recall": {
    "phase": "recall",
    "strokes": [...],
    "total_time_ms": 38000,
    "stroke_count": 4
  }
}
\end{lstlisting}

\newpage
\section{功能模块设计}

\subsection{被试信息登记}

\begin{itemize}[leftmargin=*]
  \item \textbf{输入字段}：被试编号（默认 SUBJ001）、年龄、性别（男/女/其他）、利手（右/左/双）
  \item \textbf{校验}：编号与年龄必填非空
  \item \textbf{时间戳生成}：格式 \texttt{YYYYMMDD\_HHMMSS}，作为数据文件命名前缀
\end{itemize}

\subsection{临摹阶段 (Copy Phase)}

\begin{table}[h]
\centering
\begin{tabular}{ll}
\toprule
\textbf{属性} & \textbf{值} \\
\midrule
默认时长 & 600 秒（10 分钟），可在 SKILL.md 中配置 \texttt{COPY\_TIME} \\
左侧面板 & Rey 标准刺激图（ReyFigureWidget 渲染） \\
右侧画布 & DrawingCanvas（白底画布，支持多种笔触） \\
工具栏   & 细笔 / 粗笔 / 橡皮擦 / 撤销 / 清空 / 完成绘图 \\
实时状态 & 倒计时显示（分辨率 1 秒，$<$60 秒红色警告） \\
数据捕获 & 每次 mousePress→mouseMove→mouseRelease 记录完整点序列 \\
\bottomrule
\end{tabular}
\caption{临摹阶段配置}
\end{table}

\subsubsection{画布交互细节}

\begin{itemize}[leftmargin=*]
  \item \textbf{鼠标按下}：记录起始点坐标 (x, y) 与时间戳
  \item \textbf{鼠标移动}：以固定采样率（约 60 Hz）追加路径点
  \item \textbf{鼠标释放}：将该笔画推入 strokes 数组，触发状态更新
  \item \textbf{撤销操作}：从 strokes 数组中移除最后一笔，重绘画布
  \item \textbf{橡皮擦}：切换为白色粗笔模式（逻辑擦除，非物理擦除）
\end{itemize}

\subsection{干扰任务 (Distractor Task)}

\begin{table}[h]
\centering
\begin{tabular}{ll}
\toprule
\textbf{属性} & \textbf{值} \\
\midrule
默认时长 & 60 秒，可在 SKILL.md 中配置 \texttt{DISTRACT\_TIME} \\
功能     & 全屏倒计时页面，防止被试复述图形信息 \\
显示内容 & 大字号倒计时数字 + "请稍候" 提示文字 \\
结束后   & 自动跳转回忆阶段 \\
\bottomrule
\end{tabular}
\caption{干扰任务配置}
\end{table}

\subsection{回忆阶段 (Recall Phase)}

\begin{itemize}[leftmargin=*]
  \item 与临摹阶段共享同一 ExperimentWidget，通过 \texttt{phase} 参数切换
  \item 左侧面板\textbf{不显示}刺激图（隐藏 ReyFigureWidget）
  \item 右侧画布与工具栏保持一致
  \item 默认时长 600 秒，可配置 \texttt{RECALL\_TIME}
\end{itemize}

\subsection{完成报告}

测验完成后弹出 QMessageBox 汇总：
\begin{itemize}[leftmargin=*]
  \item 被试编号
  \item 临摹笔画数
  \item 回忆笔画数
  \item 数据保存路径
\end{itemize}

同时自动保存：
\begin{itemize}[leftmargin=*]
  \item 结构化数据 → \texttt{rocf\_\{id\}\_\{timestamp\}.json}
  \item 临摹截图 → \texttt{rocf\_\{id\}\_copy\_\{timestamp\}.png}
  \item 回忆截图 → \texttt{rocf\_\{id\}\_recall\_\{timestamp\}.png}
\end{itemize}

\subsection{历史记录}

\begin{itemize}[leftmargin=*]
  \item 扫描 \texttt{OUTPUT\_DIR} 中所有 \texttt{.json} 文件
  \item 以表格展示：编号 | 性别 | 年龄 | 利手 | 日期 | 临摹笔数 | 回忆笔数 | 数据文件
  \item 支持点击表头排序
  \item 返回按钮回到主菜单
\end{itemize}

\section{评分系统映射}

\subsection{Osterrieth 18 单元体系（当前实现目标）}

参考 Rey 图形标准分解（Osterrieth, 1944），18 个单元对应图形中的关键结构元素。电子化系统通过分析绘制的几何特征判定每单元得分：

\begin{table}[h]
\centering
\begin{tabular}{cllp{4cm}}
\toprule
\textbf{单元} & \textbf{图形元素} & \textbf{判定规则} & \textbf{分值} \\
\midrule
1  & 大矩形外框   & 检测封闭四边形，边长比接近 2:1  & 0/1/2 \\
2  & 水平中线     & 检测 Y$\approx$0.5H 处接近水平的线段 & 0/1/2 \\
3  & 垂直中线     & 检测 X$\approx$0.5W 处接近垂直的线段 & 0/1/2 \\
4  & 左上对角线   & 检测从矩形左上到中心的线段  & 0/1/2 \\
5  & 右上对角线   & 检测从矩形右上到中心的线段  & 0/1/2 \\
6--18 & （其余 13 个细节单元） & — & 0/1/2 \\
\bottomrule
\end{tabular}
\caption{Osterrieth 18 单元评分体系}
\end{table}

\subsection{评分扩展计划}

\begin{table}[h]
\centering
\begin{tabular}{lll}
\toprule
\textbf{阶段} & \textbf{评分系统} & \textbf{状态} \\
\midrule
v2.0 & Osterrieth 18 单元（手动参考） & 已实现数据捕获 \\
v2.1 & 基于规则的半自动评分（参考 Sangiovanni 等, 2020） & 规划中 \\
v3.0 & BQSS 17 项定性维度支持 & 远期 \\
v4.0 & CNN 自动评分集成（参考 Langer 等, 2022） & 远期 \\
\bottomrule
\end{tabular}
\caption{评分系统扩展路线图}
\end{table}

\section{技术实现细节}

\subsection{画布渲染引擎}

\begin{lstlisting}[language=Python]
class DrawingCanvas(QWidget):
    """基于 QPainter 的自由绘画画布"""
    
    def paintEvent(self, event):
        p = QPainter(self)
        p.setRenderHint(QPainter.Antialiasing)
        # 逐笔重放所有 strokes
        for stroke in self.strokes:
            pen = QPen(QColor(stroke["color"]), stroke["width"])
            pen.setCapStyle(Qt.RoundCap)
            pen.setJoinStyle(Qt.RoundJoin)
            p.setPen(pen)
            path = QPainterPath()
            path.moveTo(stroke["points"][0]["x"], stroke["points"][0]["y"])
            for pt in stroke["points"][1:]:
                path.lineTo(pt["x"], pt["y"])
            p.drawPath(path)
\end{lstlisting}

\subsection{跨平台字体方案}

\begin{lstlisting}[language=Python]
import platform
_SYSTEM = platform.system()
if _SYSTEM == "Darwin":
    FONT_FAMILY = "PingFang SC"
elif _SYSTEM == "Windows":
    FONT_FAMILY = "Microsoft YaHei"
else:
    FONT_FAMILY = "Noto Sans CJK SC"
\end{lstlisting}

\subsection{数据目录}

\begin{lstlisting}[language=Python]
# 开发环境：项目根目录下的 output/
OUTPUT_DIR = os.path.join(os.path.dirname(__file__), "..", "output")
# 打包环境：~/Documents/ROCF测验数据/
# （由 py2app 打包时自动切换）
\end{lstlisting}

\subsection{可配置参数}

通过 Marvis SKILL.md 暴露的核心配置项：

\begin{table}[h]
\centering
\begin{tabular}{lll}
\toprule
\textbf{参数} & \textbf{默认值} & \textbf{说明} \\
\midrule
\texttt{COPY\_TIME}     & 600  & 临摹阶段时长（秒） \\
\texttt{RECALL\_TIME}   & 600  & 回忆阶段时长（秒） \\
\texttt{DISTRACT\_TIME} & 60   & 干扰任务时长（秒） \\
\texttt{WINDOW\_W}      & 1400 & 窗口宽度（像素） \\
\texttt{WINDOW\_H}      & 900  & 窗口高度（像素） \\
\bottomrule
\end{tabular}
\caption{可配置参数}
\end{table}

\section{临床应用映射}

根据文献综述中识别的各疾病 ROCF 特征模式，系统设计如下对应：

\begin{table}[h]
\centering
\begin{tabular}{p{3.8cm} p{5cm} p{5cm}}
\toprule
\textbf{疾病} & \textbf{系统捕获指标} & \textbf{临床意义} \\
\midrule
阿尔茨海默病（早期） & 临摹保留 + 延迟回忆$\downarrow$ + 再认$\downarrow$ & 海马储存缺陷，数字生物标志物 \\
TBI & 临摹$\downarrow$ + 笔序混乱 + 组织性$\downarrow$ & 执行规划损伤，可量化笔序离散度 \\
帕金森病 & 临摹轻度$\downarrow$ + 回忆$\downarrow$ + 再认保留 & 提取缺陷，DSS-ROCF 组织评分敏感 \\
ADHD & 冲动性笔触 + 遗漏细节单元 & 笔速方差大、完成时间短 \\
血管性痴呆 & 临摹显著$\downarrow$ + 碎片化组织 & BQSS 定性维度鉴别 \\
\bottomrule
\end{tabular}
\caption{临床应用映射}
\end{table}

\subsection{组织策略自动判定}

系统通过分析临摹阶段笔序自动判定策略类型（参考 Savage 等, 2000）：

\begin{itemize}[leftmargin=*]
  \item \textbf{整体式}：首笔为大矩形外框 → \texttt{strategy: "configural"}
  \item \textbf{零散式}：从独立细节开始 → \texttt{strategy: "piecemeal"}
  \item \textbf{中间型}：混合策略 → \texttt{strategy: "intermediate"}
\end{itemize}

\section{部署与分发}

\subsection{macOS .app 打包}

\begin{lstlisting}[language=bash]
python setup.py py2app
# 输出：dist/ROCF测验.app（独立应用包，无需安装 Python）
\end{lstlisting}

\subsection{DMG 分发}

\begin{lstlisting}[language=bash]
hdiutil create -volname "ROCF测验" -srcfolder dist/ROCF测验.app \
  -ov -format UDZO ROCF测验.dmg
\end{lstlisting}

\subsection{Marvis Skill 分发}

通过 GitHub Release 和 \texttt{npx skills add} 一键安装：

\begin{lstlisting}[language=bash]
npx skills add https://github.com/hyde357/rocf-test -g -y
\end{lstlisting}

\subsection{GitHub Release}

\begin{itemize}[leftmargin=*]
  \item 仓库：\url{https://github.com/hyde357/rocf-test}
  \item Release v1.0.0：包含 DMG 安装包、验证报告、18 项测试全部通过
\end{itemize}

\section{测试验证}

\subsection{测试套件}

\texttt{tests/test\_rocf.py} 包含 18 项单元测试，覆盖：

\begin{table}[h]
\centering
\begin{tabular}{lll}
\toprule
\textbf{测试类} & \textbf{条目数} & \textbf{覆盖范围} \\
\midrule
字体检测   & 3 & 跨平台字体回退链 \\
配置常量   & 4 & WINDOW\_W/H, COPY\_TIME 等 \\
数据目录   & 2 & 创建权限、读写验证 \\
PySide6 环境 & 2 & 模块可用性、版本检查 \\
应用启动   & 1 & ROCFMainWindow 初始化 \\
画布逻辑   & 3 & 笔画记录、撤销、清空 \\
数据持久化 & 2 & JSON 序列化、截图保存 \\
安装脚本   & 1 & install\_deps.py 语法 \\
\bottomrule
\end{tabular}
\caption{测试覆盖范围}
\end{table}

\subsection{测试结果}

\begin{verbatim}
Ran 18 tests in 0.074s
OK (全部通过)
\end{verbatim}

% ============================================================================
% 参考文献
% ============================================================================
\section*{附录 A：参考文献}
\addcontentsline{toc}{section}{附录 A：参考文献}

\begin{enumerate}[leftmargin=*, label={[\arabic*]}]
  \item Rey, A. (1941). L'examen psychologique dans les cas d'encéphalopathie traumatique. \textit{Archives de Psychologie}, 28, 286–340.
  \item Osterrieth, P. A. (1944). Le test de copie d'une figure complexe. \textit{Archives de Psychologie}, 30, 206–356.
  \item Meyers, J. E., \& Meyers, K. R. (1995). \textit{Rey Complex Figure Test and Recognition Trial}. PAR.
  \item Stern, R. A., et al. (1999). The Boston Qualitative Scoring System for the Rey-Osterrieth Complex Figure. \textit{Journal of Clinical and Experimental Neuropsychology}, 21(3), 401–417.
  \item Savage, C. R., et al. (2000). Organizational strategies mediate nonverbal memory impairment in obsessive-compulsive disorder. \textit{Biological Psychiatry}, 45(7), 905–916.
  \item Langer, N., et al. (2022). Automated scoring of the Rey-Osterrieth Complex Figure Test using deep learning. \textit{Neuropsychology}, 36(5), 456–467.
  \item Zhang, L., et al. (2018). Normative data for the ROCF in a Chinese population. \textit{Applied Neuropsychology: Adult}, 25(3), 233–241.
  \item Zhang, Y., et al. (2021). Digital pen-based administration of the ROCF captures latent neuropsychological structure. \textit{Assessment}, 28(6), 1562–1574.
\end{enumerate}

\section*{附录 B：项目文件结构}
\addcontentsline{toc}{section}{附录 B：项目文件结构}

\begin{verbatim}
rocf-test/
├── SKILL.md                  # Marvis Skill 描述与配置
├── README.md                 # 项目说明
├── .gitignore
├── assets/
│   └── rocf_qt.py            # 主程序源码（831 行）
├── scripts/
│   └── install_deps.py       # 依赖安装脚本
├── tests/
│   └── test_rocf.py          # 测试套件（18 项）
├── VALIDATION_REPORT.md      # 验证报告
└── ROCF测验.dmg              # macOS 安装包（36 MB）
\end{verbatim}

\end{document}
